% !TEX root = ../main.tex
\chapter{Suivi d'Exp\'eriences --- MLflow, Weights \& Biases}
\label{ch:suivi-experiences}

\section{Introduction}

Entra\^iner un mod\`ele de Machine Learning implique de tester des dizaines,
voire des centaines, de combinaisons d'hyperparam\`etres, d'architectures et
de jeux de donn\'ees. Sans un syst\`eme de suivi structur\'e, les r\'esultats
se perdent et la reproductibilit\'e devient illusoire.

\begin{antipattern}[title=\textbf{L'anti-pattern du tableur}]
Beaucoup de data scientists commencent par noter leurs r\'esultats dans un
fichier Excel ou Google Sheets~:
\begin{itemize}
  \item Les entr\'ees sont incompl\`etes (on oublie de noter un param\`etre).
  \item Aucun lien avec le code ou les donn\'ees utilis\'es.
  \item Impossible de reproduire une exp\'erience pass\'ee.
  \item Le fichier devient vite illisible avec des centaines de lignes.
\end{itemize}
Les outils de suivi d'exp\'eriences r\'esolvent tous ces probl\`emes.
\end{antipattern}

\begin{definition}[Suivi d'exp\'eriences]
Le \textbf{suivi d'exp\'eriences} (\emph{experiment tracking}) consiste \`a
enregistrer automatiquement, pour chaque ex\'ecution d'entra\^inement~:
les \textbf{param\`etres}, les \textbf{m\'etriques}, les \textbf{artefacts}
(mod\`eles, graphiques), le \textbf{code source} et l'\textbf{environnement}.
\end{definition}

%% ============================================================
\section{MLflow}
\label{sec:mlflow}
%% ============================================================

\index{MLflow}

\subsection{Pr\'esentation}

\pkg{MLflow} est une plateforme open-source cr\'e\'ee par Databricks pour
g\'erer le cycle de vie complet du ML. Elle comprend quatre composants~:

\begin{formules}[title=\textbf{Composants de MLflow}]
\begin{description}
  \item[MLflow Tracking] Enregistrement des param\`etres, m\'etriques et artefacts.
  \item[MLflow Projects] Packaging reproductible du code ML.
  \item[MLflow Models] Format standard pour le d\'eploiement de mod\`eles.
  \item[MLflow Model Registry] Gestion du cycle de vie des mod\`eles
    (staging, production, archiv\'e).
\end{description}
\end{formules}

\subsection{Installation et d\'emarrage}

\begin{codeblock}[title=\textbf{Installation de MLflow}]
\begin{minted}{bash}
pip install mlflow

# Lancer l'interface web
mlflow ui --port 5000
# Accessible sur http://localhost:5000
\end{minted}
\end{codeblock}

\subsection{API de suivi (Tracking API)}
\index{MLflow!Tracking API}

\begin{codeblock}[title=\textbf{Concepts fondamentaux de MLflow Tracking}]
\begin{minted}{python}
import mlflow

# Cr\'eer ou s\'electionner une exp\'erience
mlflow.set_experiment("classification-vin")

# D\'emarrer un run
with mlflow.start_run(run_name="random-forest-v1"):
    # Enregistrer des param\`etres
    mlflow.log_param("n_estimators", 100)
    mlflow.log_param("max_depth", 10)
    mlflow.log_param("random_state", 42)

    # Enregistrer des m\'etriques
    mlflow.log_metric("accuracy", 0.94)
    mlflow.log_metric("f1_score", 0.93)

    # Enregistrer des m\'etriques au fil des \'epoques
    for epoch in range(50):
        mlflow.log_metric("loss", loss_val, step=epoch)

    # Enregistrer des artefacts
    mlflow.log_artifact("plots/confusion_matrix.png")
    mlflow.log_artifact("models/model.pkl")
\end{minted}
\end{codeblock}

\subsection{Entra\^inement complet avec MLflow}

\begin{codeblock}[title=\textbf{Exemple complet~: classification avec MLflow}]
\begin{minted}{python}
import mlflow
import mlflow.sklearn
from sklearn.datasets import load_wine
from sklearn.model_selection import train_test_split
from sklearn.ensemble import RandomForestClassifier
from sklearn.metrics import accuracy_score, f1_score
import matplotlib.pyplot as plt
from sklearn.metrics import ConfusionMatrixDisplay

# Charger les donn\'ees
X, y = load_wine(return_X_y=True)
X_train, X_test, y_train, y_test = train_test_split(
    X, y, test_size=0.2, random_state=42
)

# Configuration
params = {
    "n_estimators": 200,
    "max_depth": 15,
    "min_samples_split": 5,
    "random_state": 42,
}

mlflow.set_experiment("wine-classification")

with mlflow.start_run(run_name="rf-optimized"):
    # Enregistrer les param\`etres
    mlflow.log_params(params)

    # Entra\^iner le mod\`ele
    model = RandomForestClassifier(**params)
    model.fit(X_train, y_train)

    # Pr\'edire et \'evaluer
    y_pred = model.predict(X_test)
    acc = accuracy_score(y_test, y_pred)
    f1 = f1_score(y_test, y_pred, average="weighted")

    mlflow.log_metric("accuracy", acc)
    mlflow.log_metric("f1_weighted", f1)

    # Matrice de confusion
    fig, ax = plt.subplots()
    ConfusionMatrixDisplay.from_predictions(
        y_test, y_pred, ax=ax
    )
    fig.savefig("confusion_matrix.png")
    mlflow.log_artifact("confusion_matrix.png")

    # Enregistrer le mod\`ele
    mlflow.sklearn.log_model(model, "random-forest-model")

    print(f"Accuracy: {acc:.4f}, F1: {f1:.4f}")
\end{minted}
\end{codeblock}

\subsection{MLflow Model Registry}
\index{MLflow!Model Registry}

\begin{codeblock}[title=\textbf{Utilisation du Model Registry}]
\begin{minted}{python}
import mlflow

# Enregistrer un mod\`ele dans le registry
model_uri = "runs:/<run_id>/random-forest-model"
mlflow.register_model(model_uri, "WineClassifier")

# Charger un mod\`ele depuis le registry
from mlflow import MlflowClient

client = MlflowClient()

# Promouvoir en staging
client.transition_model_version_stage(
    name="WineClassifier",
    version=1,
    stage="Staging"
)

# Promouvoir en production
client.transition_model_version_stage(
    name="WineClassifier",
    version=1,
    stage="Production"
)

# Charger le mod\`ele en production
model = mlflow.pyfunc.load_model(
    "models:/WineClassifier/Production"
)
\end{minted}
\end{codeblock}

\begin{bonnepratique}[title=\textbf{Workflow du Model Registry}]
\begin{enumerate}
  \item \textbf{None}~: le mod\`ele vient d'\^etre enregistr\'e.
  \item \textbf{Staging}~: le mod\`ele est en cours de validation.
  \item \textbf{Production}~: le mod\`ele est d\'eploy\'e en production.
  \item \textbf{Archived}~: le mod\`ele est retir\'e mais conserv\'e.
\end{enumerate}
\end{bonnepratique}

%% ============================================================
\section{Weights \& Biases (W\&B)}
\label{sec:wandb}
%% ============================================================

\index{Weights \& Biases}

\subsection{Pr\'esentation}

\pkg{Weights \& Biases} (W\&B) est une plateforme cloud de suivi
d'exp\'eriences offrant une interface web riche, des visualisations
interactives et des fonctionnalit\'es de collaboration avanc\'ees.

\begin{codeblock}[title=\textbf{Installation et authentification}]
\begin{minted}{bash}
pip install wandb
wandb login  # Saisir la cl\'e API depuis wandb.ai
\end{minted}
\end{codeblock}

\subsection{API de suivi W\&B}
\index{Weights \& Biases!API}

\begin{codeblock}[title=\textbf{Entra\^inement complet avec W\&B}]
\begin{minted}{python}
import wandb
from sklearn.datasets import load_wine
from sklearn.model_selection import train_test_split
from sklearn.ensemble import GradientBoostingClassifier
from sklearn.metrics import accuracy_score, f1_score

# Initialiser W&B
wandb.init(
    project="wine-classification",
    name="gradient-boosting-v1",
    config={
        "n_estimators": 150,
        "learning_rate": 0.1,
        "max_depth": 5,
        "random_state": 42,
    }
)
config = wandb.config

# Charger les donn\'ees
X, y = load_wine(return_X_y=True)
X_train, X_test, y_train, y_test = train_test_split(
    X, y, test_size=0.2, random_state=config.random_state
)

# Entra\^iner le mod\`ele
model = GradientBoostingClassifier(
    n_estimators=config.n_estimators,
    learning_rate=config.learning_rate,
    max_depth=config.max_depth,
    random_state=config.random_state,
)

# Suivi \'epoque par \'epoque (staged_predict)
model.fit(X_train, y_train)
for i, y_pred_staged in enumerate(
    model.staged_predict(X_test)
):
    acc = accuracy_score(y_test, y_pred_staged)
    wandb.log({"epoch": i, "accuracy": acc})

# M\'etriques finales
y_pred = model.predict(X_test)
final_acc = accuracy_score(y_test, y_pred)
final_f1 = f1_score(y_test, y_pred, average="weighted")

wandb.log({
    "final_accuracy": final_acc,
    "final_f1": final_f1,
})

# Enregistrer le mod\`ele comme artefact
artifact = wandb.Artifact("wine-model", type="model")
artifact.add_file("models/model.pkl")
wandb.log_artifact(artifact)

wandb.finish()
\end{minted}
\end{codeblock}

\subsection{Sweeps~: recherche d'hyperparam\`etres}
\index{Weights \& Biases!Sweeps}

\begin{codeblock}[title=\textbf{Configuration d'un Sweep W\&B}]
\begin{minted}{yaml}
# sweep_config.yaml
program: train.py
method: bayes  # random, grid, bayes
metric:
  name: final_accuracy
  goal: maximize
parameters:
  n_estimators:
    values: [50, 100, 200, 500]
  learning_rate:
    min: 0.001
    max: 0.3
  max_depth:
    values: [3, 5, 7, 10]
\end{minted}
\end{codeblock}

\begin{codeblock}[title=\textbf{Lancer un Sweep}]
\begin{minted}{bash}
# Cr\'eer le sweep
wandb sweep sweep_config.yaml
# Retourne un sweep_id

# Lancer les agents
wandb agent <username>/<project>/<sweep_id>
\end{minted}
\end{codeblock}

%% ============================================================
\section{Hydra pour la gestion de configuration}
\label{sec:hydra}
%% ============================================================

\index{Hydra}

\pkg{Hydra} est un framework de configuration d\'evelopp\'e par Meta qui
simplifie la gestion des param\`etres dans les projets ML.

\begin{codeblock}[title=\textbf{Structure de configuration Hydra}]
\begin{minted}{yaml}
# config/config.yaml
defaults:
  - model: random_forest
  - data: wine

training:
  test_size: 0.2
  random_state: 42

# config/model/random_forest.yaml
name: RandomForest
n_estimators: 100
max_depth: 10

# config/model/gradient_boosting.yaml
name: GradientBoosting
n_estimators: 200
learning_rate: 0.1

# config/data/wine.yaml
name: wine
path: data/raw/wine.csv
\end{minted}
\end{codeblock}

\begin{codeblock}[title=\textbf{Utilisation de Hydra avec MLflow}]
\begin{minted}{python}
import hydra
from omegaconf import DictConfig
import mlflow

@hydra.main(config_path="config", config_name="config",
            version_base=None)
def train(cfg: DictConfig):
    mlflow.set_experiment(cfg.data.name)

    with mlflow.start_run():
        # Les param\`etres viennent de la config Hydra
        mlflow.log_params({
            "model": cfg.model.name,
            "n_estimators": cfg.model.n_estimators,
            "test_size": cfg.training.test_size,
        })

        # ... entra\^inement ...

        mlflow.log_metric("accuracy", accuracy)

if __name__ == "__main__":
    train()
\end{minted}
\end{codeblock}

\begin{codeblock}[title=\textbf{Override de configuration en ligne de commande}]
\begin{minted}{bash}
# Utiliser la config par d\'efaut
python train.py

# Changer de mod\`ele
python train.py model=gradient_boosting

# Override un param\`etre
python train.py model.n_estimators=500

# Multi-run (grid search)
python train.py --multirun \
    model=random_forest,gradient_boosting \
    model.n_estimators=100,200,500
\end{minted}
\end{codeblock}

%% ============================================================
\section{Architecture du suivi d'exp\'eriences}
\label{sec:tracking-architecture}
%% ============================================================

\begin{figure}[htbp]
\centering
\begin{tikzpicture}[
    node distance=1.5cm and 2cm,
    block/.style={rectangle, draw, rounded corners, minimum width=2.8cm,
                  minimum height=1cm, align=center, font=\small},
    server/.style={rectangle, draw, rounded corners, minimum width=3cm,
                   minimum height=1.2cm, align=center, font=\small,
                   double, double distance=1.5pt},
    storage/.style={cylinder, draw, shape border rotate=90,
                    minimum width=1.8cm, minimum height=1.2cm,
                    align=center, font=\small},
    arrow/.style={-{Stealth}, thick},
  ]

  % Sources
  \node[block, fill=blue!15] (script) {Script\\d'entra\^inement};
  \node[block, fill=blue!15, below=1cm of script] (notebook) {Notebook\\Jupyter};
  \node[block, fill=blue!15, below=1cm of notebook] (hydra) {Hydra\\Config};

  % Serveur de tracking
  \node[server, fill=green!20, right=3cm of notebook] (tracker) {Serveur de\\suivi\\(MLflow / W\&B)};

  % Stockage
  \node[storage, fill=orange!15, right=2.5cm of tracker] (db) {Base de\\donn\'ees\\m\'etriques};
  \node[storage, fill=yellow!15, above=1cm of db] (artifacts) {Stockage\\artefacts};
  \node[block, fill=purple!15, below=1cm of db] (ui) {Interface\\Web (UI)};

  % Fl\`eches
  \draw[arrow] (script) -- (tracker);
  \draw[arrow] (notebook) -- (tracker);
  \draw[arrow] (hydra) -- (tracker);
  \draw[arrow] (tracker) -- (db);
  \draw[arrow] (tracker) -- (artifacts);
  \draw[arrow] (db) -- (ui);
  \draw[arrow] (artifacts) -- (ui);

  % \'Etiquettes
  \node[font=\scriptsize, above=0.1cm of script.east, anchor=south west]
    {param\`etres, m\'etriques};
  \node[font=\scriptsize, below right=0.1cm of tracker.east, anchor=north west]
    {SQL / API};

\end{tikzpicture}
\caption{Architecture g\'en\'erale d'un syst\`eme de suivi d'exp\'eriences.
Les scripts et notebooks envoient les donn\'ees au serveur de suivi, qui les
persiste dans une base de donn\'ees et un stockage d'artefacts, accessibles
via une interface web.}
\label{fig:tracking-architecture}
\end{figure}

%% ============================================================
\section{Comparaison des outils de suivi}
\label{sec:tracking-tools-comparison}
%% ============================================================

\index{Neptune}\index{ClearML}

\begin{table}[htbp]
\centering
\caption{Comparaison des outils de suivi d'exp\'eriences}
\label{tab:tracking-tools}
\small
\begin{tabular}{lcccc}
\toprule
\textbf{Crit\`ere} & \textbf{MLflow} & \textbf{W\&B} & \textbf{Neptune} & \textbf{ClearML} \\
\midrule
Licence & Open-source & Freemium & Freemium & Open-source \\
H\'ebergement & Self-hosted & Cloud & Cloud & Self/Cloud \\
Interface web & Bonne & Excellente & Excellente & Bonne \\
Model Registry & Oui & Oui & Oui & Oui \\
Sweeps/HPO & Non natif & Oui & Oui & Oui \\
Collaboration & Basique & Avanc\'ee & Avanc\'ee & Bonne \\
Int\'egration DL & Bonne & Excellente & Excellente & Bonne \\
Taille \'equipe & Toutes & Petite--grande & Moyenne--grande & Toutes \\
Co\^ut production & Gratuit & Payant & Payant & Gratuit \\
\bottomrule
\end{tabular}
\end{table}

\begin{remarque}
\textbf{MLflow} est le choix privil\'egi\'e pour les \'equipes souhaitant garder
le contr\^ole total de leur infrastructure (self-hosted). \textbf{W\&B} offre
la meilleure exp\'erience utilisateur pour les \'equipes de recherche qui
pr\'ef\`erent une solution cloud cl\'e en main.
\end{remarque}

%% ============================================================
\section{Mini-projet~: suivi d'exp\'eriences pour le projet fil rouge}
\label{sec:miniprojet-tracking}
%% ============================================================

\begin{miniprojet}[title=\textbf{Instrumenter le projet des chapitres 1--3 avec MLflow}]
\textbf{Objectif}~: Ajouter le suivi d'exp\'eriences au pipeline de
classification cr\'e\'e dans les chapitres pr\'ec\'edents.

\begin{enumerate}
  \item Installer et lancer MLflow~:
\begin{minted}{bash}
pip install mlflow
mlflow ui --port 5000 &
\end{minted}

  \item Modifier le script d'entra\^inement pour int\'egrer MLflow~:
\begin{minted}{python}
import mlflow
import mlflow.sklearn
from sklearn.ensemble import RandomForestClassifier
from sklearn.metrics import accuracy_score, f1_score
import json

mlflow.set_experiment("projet-fil-rouge")

with mlflow.start_run(run_name="baseline-rf"):
    params = {
        "n_estimators": 100,
        "max_depth": 10,
        "random_state": 42,
    }
    mlflow.log_params(params)

    model = RandomForestClassifier(**params)
    model.fit(X_train, y_train)

    y_pred = model.predict(X_test)
    metrics = {
        "accuracy": accuracy_score(y_test, y_pred),
        "f1": f1_score(y_test, y_pred, average="weighted"),
    }
    mlflow.log_metrics(metrics)

    # Sauvegarder les m\'etriques pour DVC
    with open("metrics/scores.json", "w") as f:
        json.dump(metrics, f, indent=2)

    mlflow.sklearn.log_model(model, "model")
\end{minted}

  \item Cr\'eer une configuration Hydra pour varier les hyperparam\`etres.

  \item Lancer au moins 5 exp\'eriences avec des param\`etres diff\'erents.

  \item Utiliser l'interface MLflow pour comparer les runs et
        identifier la meilleure configuration.

  \item Enregistrer le meilleur mod\`ele dans le Model Registry et le
        promouvoir en \texttt{Staging}.
\end{enumerate}
\end{miniprojet}

%% ============================================================
\section{Exercices}
\label{sec:exos-suivi-experiences}
%% ============================================================

\begin{exercice}[Prise en main de MLflow]
\textbf{1/3}
Installez MLflow et enregistrez une exp\'erience simple~:
\begin{enumerate}
  \item Entra\^inez un mod\`ele \texttt{LogisticRegression} sur le dataset Iris.
  \item Enregistrez les param\`etres \texttt{C}, \texttt{solver} et
        \texttt{max\_iter}.
  \item Enregistrez les m\'etriques \texttt{accuracy} et \texttt{f1\_macro}.
  \item Lancez l'interface MLflow et v\'erifiez que le run appara\^it.
  \item Relancez avec des param\`etres diff\'erents et comparez dans l'UI.
\end{enumerate}
\end{exercice}

\begin{exercice}[Sweeps W\&B et analyse]
\textbf{2/3}
Utilisez W\&B Sweeps pour optimiser les hyperparam\`etres d'un mod\`ele
\texttt{GradientBoosting}~:
\begin{enumerate}
  \item Cr\'eez un fichier \file{sweep\_config.yaml} explorant~:
    \begin{itemize}
      \item \texttt{n\_estimators}~: 50, 100, 200, 500.
      \item \texttt{learning\_rate}~: entre 0.001 et 0.5 (\'echelle log).
      \item \texttt{max\_depth}~: 3, 5, 7, 10.
    \end{itemize}
  \item Lancez un sweep bay\'esien avec 20 runs.
  \item Analysez les r\'esultats dans l'interface W\&B~: parallel coordinates,
        importance des hyperparam\`etres.
  \item R\'edigez un rapport W\&B r\'esumant vos conclusions.
\end{enumerate}
\end{exercice}

\begin{exercice}[Pipeline complet Hydra + MLflow + DVC]
\textbf{3/3}
Construisez un pipeline ML complet int\'egrant les trois outils~:
\begin{enumerate}
  \item Utilisez \textbf{DVC} pour versionner les donn\'ees et d\'efinir les
        \'etapes du pipeline (\file{dvc.yaml}).
  \item Utilisez \textbf{Hydra} pour g\'erer les configurations~: cr\'eez au
        moins trois fichiers de config (un par mod\`ele).
  \item Utilisez \textbf{MLflow} pour suivre chaque exp\'erience
        automatiquement.
  \item Int\'egrez le tout pour que \cli{dvc repro} ex\'ecute le pipeline et
        enregistre les r\'esultats dans MLflow.
  \item Comparez les r\'esultats de trois architectures diff\'erentes et
        enregistrez la meilleure dans le Model Registry.
\end{enumerate}
\end{exercice}

%% ============================================================
\section{Aide-m\'emoire}
\label{sec:cheatsheet-suivi-experiences}
%% ============================================================

\begin{cheatsheet}[title=\textbf{Suivi d'exp\'eriences --- Aide-m\'emoire}]
\small
\begin{tabular}{p{4.5cm}p{8.5cm}}
\toprule
\textbf{Action} & \textbf{Commande / Code} \\
\midrule
\multicolumn{2}{l}{\textbf{MLflow}} \\
\midrule
Lancer l'UI & \cli{mlflow ui --port 5000} \\
Cr\'eer une exp\'erience & \texttt{mlflow.set\_experiment("nom")} \\
D\'emarrer un run & \texttt{mlflow.start\_run(run\_name="nom")} \\
Logger param\`etres & \texttt{mlflow.log\_params(dict)} \\
Logger m\'etriques & \texttt{mlflow.log\_metric("nom", val, step)} \\
Logger un mod\`ele & \texttt{mlflow.sklearn.log\_model(m, "nom")} \\
Enregistrer au registry & \texttt{mlflow.register\_model(uri, "nom")} \\
\midrule
\multicolumn{2}{l}{\textbf{Weights \& Biases}} \\
\midrule
Connexion & \cli{wandb login} \\
Initialiser & \texttt{wandb.init(project, name, config)} \\
Logger m\'etriques & \texttt{wandb.log(\{"metric": val\})} \\
Logger artefact & \texttt{wandb.log\_artifact(artifact)} \\
Cr\'eer un sweep & \cli{wandb sweep config.yaml} \\
Lancer un agent & \cli{wandb agent user/proj/sweep\_id} \\
Terminer & \texttt{wandb.finish()} \\
\midrule
\multicolumn{2}{l}{\textbf{Hydra}} \\
\midrule
D\'ecorateur principal & \texttt{@hydra.main(config\_path, config\_name)} \\
Override param\`etre & \cli{python train.py model.lr=0.01} \\
Multi-run & \cli{python train.py --multirun model=a,b,c} \\
\bottomrule
\end{tabular}
\end{cheatsheet}
